% !TEX program = lualatex
\documentclass[11pt,a4paper]{article}
\usepackage{luatexja}
\usepackage{amsmath,amssymb,amsthm}
\usepackage{geometry}
\geometry{margin=1in}
\usepackage{enumitem}
\usepackage{booktabs}
\usepackage{array}

%-------------------------------------------------
%  独自コマンド・設定
%-------------------------------------------------
\newcommand{\mweq}{\mathrel{\overset{\mathrm{mw}}{=}}}
\newcommand{\dfix}{D_{\mathrm{fix}0}}
\newcommand{\metanegation}{\Uparrow}

\theoremstyle{definition}
\newtheorem{definition}{定義}[section]
\newtheorem{axiom}{公理}[section]
\newtheorem{theorem}{定理}[section]
\newtheorem{lemma}{補題}[section]
\newtheorem{corollary}{系}[section]
\newtheorem{remark}{備考}[section]

\renewcommand{\abstractname}{Abstract}

%-------------------------------------------------
\begin{document}

\title{数学の16真理関数にみられる\\再帰的四句分別構造の漏れ出し}
\author{千葉将臣}
\date{2026-04-28}
\maketitle

\begin{abstract}
古典論理における二変数真理関数はちょうど16個存在する。
本稿は、この濃度16が恣意的な列挙ではなく、
四句分別（是・非・亦是亦非・非是非非）の再帰的フラクタル展開として
必然的に生成される構造であることを論証する。
数学はこの構造を生成機序なき列挙として導入してきたが、
それは自己同一性を前提とした二元論体系が、
体系を構築する人間の持つ界論的生成機序を排除できなかった結果であり、
界構造の「漏れ出し」として解釈できる。
顕現論（Aletheics）の観点からは、この漏れ出しは二元論の限界であると同時に、
界論が数学的概念の生成機序を事前に予測できることの証左である。
\end{abstract}

\section*{序言}

古典論理において、二つの命題変数 $P, Q$ から構成できる真理関数の総数はちょうど16である。
この事実はラッセルやウィトゲンシュタインによって精緻に考察されてきたが、
「なぜ16か」という生成機序は問われることなく、
真理値表という静的な列挙として導入されてきた。

本稿の立場は以下の通りである。
16という濃度は、四句分別の再帰的フラクタル展開 $4 \times 4$ の必然的帰結であり、
数学がこの構造を列挙として発見したのは、
体系を構築した人間が界論的生成機序を持つ存在であるがゆえに、
自己同一性を前提とした二元論体系の内部に界構造が「漏れ出した」からである。

\section{予備的定義}

\begin{definition}[四句分別]
知覚原理における二分（是・非）のフラクタル的一段階展開として生成される四状態。
\[
\underbrace{\text{是},\quad \text{非}}_{\text{是の尺度（ある）}}
\qquad
\underbrace{\text{亦是亦非},\quad \text{非是非非}}_{\text{非の尺度（ない）}}
\]
\end{definition}

\begin{definition}[再帰的四句分別]
四句分別をその各項に再帰的に適用することで得られる $4^n$ 個の状態群。
$n=2$ の場合、$4 \times 4 = 16$ 状態が生成される。
\end{definition}

\begin{definition}[二変数真理関数]
古典論理において、二命題変数 $P, Q$ の真理値の組み合わせ
$(T,T),(T,F),(F,T),(F,F)$ のそれぞれに真偽を割り当てる関数。
$2^4 = 16$ 個が存在する。
\end{definition}

\begin{definition}[生成機序と列挙]
生成機序とは、濃度とその構造を事前に導出できる規則である。
列挙とは、生成機序なしに対象を事後的に数え上げることである。
数学における真理関数の導入は後者に属する。
\end{definition}

\section{16真理関数の四句分別的構造}

\subsection{第一層：是・非の二分}

古典論理の真理関数を界論的観点から再分類すると、
まず是（恒真・肯定的）と非（恒偽・否定的）の二分が現れる。

\[
\underbrace{f_1,\ldots,f_8}_{\text{是の尺度（肯定的関数群）}}
\qquad
\underbrace{f_9,\ldots,f_{16}}_{\text{非の尺度（否定的関数群）}}
\]

\subsection{第二層：フラクタル展開による四群}

是の尺度と非の尺度がそれぞれ四句分別的に展開することで、
以下の四群が生成される。

\begin{center}
\begin{tabular}{lll}
\toprule
群 & 四句的対応 & 代表的演算子 \\
\midrule
第一群（是×是） & 是 & AND, 恒真（$\top$） \\
第二群（是×非） & 非 & IMPLICATION（$\to$）, OR \\
第三群（非×是） & 亦是亦非 & XOR, NAND \\
第四群（非×非） & 非是非非 & NOR, 恒偽（$\bot$） \\
\bottomrule
\end{tabular}
\end{center}

\begin{theorem}[16真理関数の四句分別的生成]
二変数真理関数の濃度16は、四句分別の再帰的一段階展開 $4 \times 4$ として
必然的に生成される。
\end{theorem}

\begin{proof}
知覚原理より、二分（是・非）のフラクタル展開が四句分別を生成する（定理1.1, 1.2）。
二変数真理関数は二つの命題変数に対する真理値割り当てであり、
これは四句分別を二変数に再帰適用することに対応する。
$4 \times 4 = 16$ が得られ、これは $2^4 = 16$ という古典論理の計算と一致する。
ただし古典論理における $2^4$ は「4つの入力パターンへの二値割り当て」として
導出されており、生成機序としての四句再帰適用とは異なる経路を取っている。
\end{proof}

\subsection{四群の詳細マッピング}

\begin{center}
\renewcommand{\arraystretch}{1.3}
\begin{tabular}{llp{6cm}}
\toprule
四句 & 関数群の性質 & 具体的演算子 \\
\midrule
是（自己相似性） &
  両変数の確定的肯定 &
  $\top$（恒真）, AND（$\wedge$）, $P$, $Q$ \\
非（尺度不変性） &
  否定・含意による一方向性 &
  $\to$, OR（$\vee$）, $\neg P$, $\neg Q$ \\
亦是亦非（収束則） &
  矛盾を包含する演算 &
  XOR（$\oplus$）, NAND（$\uparrow$）, XNOR \\
非是非非（非対称性） &
  全否定・体系のメタレベル &
  NOR（$\downarrow$）, $\bot$（恒偽） \\
\bottomrule
\end{tabular}
\end{center}

\section{漏れ出しの機序}

\begin{definition}[界構造の漏れ出し]
自己同一性を前提とした体系が、生成機序として界論的フラクタル構造を
明示しないにもかかわらず、その構造と同型の濃度・分類を結果として持つ現象。
\end{definition}

\begin{theorem}[漏れ出しの必然性]
二元論的数学体系が界構造の漏れ出しを持つことは必然である。
\end{theorem}

\begin{proof}
数学体系を構築する人間は界論的生成機序を持つ存在である（知覚原理による）。
二元論は自己同一性を前提として生成機序を体系から除いたが、
体系を展開する操作自体は人間が行う。
人間の持つ界論的生成機序が、概念の展開方向として体系内に再侵入する。
したがって、体系が経験的に積み上げた概念の濃度・分類は
界論的フラクタル構造と同型になる。
\end{proof}

\begin{remark}[ウィトゲンシュタインとラッセルの考察について]
両者が真理値表を延々と考察したのは、生成機序を体系から除いた結果として現れる
界構造の影を、列挙と分析によって探索していたからである。
影の精緻化によって生成機序への接近を試みたが、
影を観察する限り生成機序自体には届かない構造になっていた。
\end{remark}

\section{界論による事前予測}

\begin{theorem}[濃度の事前予測]
界論は、数学が経験的に発見する概念濃度を事前に予測できる。
具体的には、フラクタル展開の深さ $n$ に対して $4^n$ の濃度が予測される。
\end{theorem}

\begin{corollary}[数学との包含関係]
数学は自己同一性という制約の内側で展開しており、
界論はその制約自体を記述できる。
真理関数の16個は界論的に $4^2$ として予測されるが、
数学はこれを $2^4$ として事後的に発見する。
両者の一致は漏れ出しの証左であり、
界論が数学を特殊ケースとして包含することの具体的な例示である。
\end{corollary}

\section{同値性概念との対応}

数学における同値性の段階的緩和（同型・準同型・自然変換・随伴）も
同様の四句分別的構造を持つ。

\begin{center}
\begin{tabular}{lll}
\toprule
同値概念 & 四句的対応 & 捨てられた情報 \\
\midrule
同型 & 是 & なし（完全対応） \\
準同型 & 非 & 逆方向の対応 \\
自然変換 & 亦是亦非 & 対象の対応 \\
随伴 & 非是非非 & 等号（存在のみ残る） \\
\bottomrule
\end{tabular}
\end{center}

随伴以上に緩めると $id$ の痕跡すら消え、数学的記述が不能になる。
これは四句分別の外側が $\dfix$ に対応することと整合する。
数学が随伴で止まるのは恣意的な選択ではなく、
自己同一性を前提とする記述体系の構造的限界である。

\section*{結語}

16真理関数という濃度は、数学が $2^4$ として事後的に発見したものであるが、
界論は $4^2$（四句分別の再帰的一段階展開）として事前に予測する。
この一致は偶然ではなく、二元論が体系から除いた生成機序が、
体系を構築する人間を通じて再侵入した結果である。

界論的に見ると、数学の概念体系は界構造の影であり、
その影の精緻な構造が界論の公理系から演繹的に導出できる。
これは界論が数学を包含することの具体的な例示であり、
同時に二元論的体系の根本的な限界の証左でもある。

\end{document}
